We represent a feasible integer-slot schedule by its completion-time vector.
Within each slot, assign jobs in increasing job order to machines in increasing
order. For a set of ready jobs,
$\operatorname{first}_k$ returns its first $k$ members in increasing order
of $M_j(\alpha)$, with ties broken by increasing job number.
An isolated zero has a neighborhood containing no other zeros. In the
return statement, $\operatorname*{arg\,min}$ selects the lexicographically
first minimizing vector, comparing vectors at their first differing coordinate.
Use a fixed deterministic polynomial-time LP algorithm to choose a rational
optimal solution in the step below.

\begin{algorithm}[H]
\DontPrintSemicolon
\caption{Deterministic common-quantile scheduling}
\label{alg:rounding}
\KwIn{$\mathcal O$, the precedence relation, $(w_j)_{j\in\mathcal O}$, $m$}
\KwOut{A feasible completion-time vector}
\If{$m\ge n$}{\KwRet{$(a_j)_{j\in\mathcal O}$}\;}
Compute an optimal LP solution $(x_{i,s})_{i\in\mathcal O,\,1\le s\le n}$\;
$\mathcal C\gets\{0,1\}\cup
  \bigl\{F_i(s):i\in\mathcal O,\ 1\le s\le n\bigr\}$\;
$\mathcal C\gets\mathcal C\cup
  \bigl\{\alpha\in(0,1):\text{$\alpha$ is an isolated zero of $M_i-M_k$ for
  some distinct $i,k\in\mathcal O$}\bigr\}$\;
$\mathcal S\gets\varnothing$\;
\ForEach{consecutive $\alpha_-<\alpha_+$ in $\mathcal C$}{
  $\alpha\gets(\alpha_-+\alpha_+)/2$\;
  $\mathcal W\gets\mathcal O$\;
  $s\gets1$\;
  \While{$\mathcal W\ne\varnothing$}{
    $\mathcal U\gets\{j\in\mathcal W:\{i\in\mathcal O:i\to j\}\cap\mathcal W=\varnothing\}$\;
    $\mathcal V\gets\operatorname{first}_{\min\{m,|\mathcal U|\}}(\mathcal U)$\;
    $(C_j)_{j\in\mathcal V}\gets(s)_{j\in\mathcal V}$\;
    $\mathcal W\gets\mathcal W\setminus\mathcal V$\;
    $s\gets s+1$\;
  }
  $\mathcal S\gets\mathcal S\cup\{(C_j)_{j\in\mathcal O}\}$\;
}
\KwRet{$\displaystyle\operatorname*{arg\,min}_{(C_j)_{j\in\mathcal O}\in\mathcal S}
                  \sum_{j\in\mathcal O}w_j\cdot C_j$}\;
\end{algorithm}

The completed jobs form a predecessor-closed set, so checking incoming arcs
suffices to identify ready jobs. The set $\mathcal U$ is nonempty whenever
$\mathcal W$ is nonempty because the precedence graph is acyclic.
If two jobs' quantile functions agree throughout an interval,
increasing job number resolves their tie and no extra breakpoint is
needed. Section~\ref{sec:foundations-derandomization} proves polynomial
running time and bounds the cost returned by Algorithm~\ref{alg:rounding}.
